\documentclass[]{article}
\usepackage[spanish,es-nodecimaldot]{babel} %load spanish as the language
\usepackage[latin1]{inputenc} %allow to put stress in words without using latex commands
%\usepackage[utf8]{inputenc} %use this in overleaf and comment the later
\usepackage{lmodern}
\usepackage[T1]{fontenc} %fix word stressing in the pdf file
\usepackage{amsmath} %Use special AMS-math environments
\usepackage{graphicx} % Enable the inclusion of images.
\usepackage {subfigure} % Enable inclusion of subfigures.
\usepackage{floatrow}
\usepackage{placeins}
\usepackage{lastpage} % Required to print the total number of pages
\usepackage{array} % More control over tables
\usepackage{tikz}
\usepackage{titlesec}
\usepackage[left=2cm,right=2cm,top=3cm,bottom=3cm]{geometry} % Adjust page margins
\usepackage{fancyhdr}

%----------------------------------------------------------------------------------------
%	MODIFY SECTION FORMAT
%----------------------------------------------------------------------------------------
%modify the format of the section headers, removes numeration and add lines
\titleformat{\section}{\normalfont\large\bfseries}{}{0pt}{\titlerule\newline\llap{}}[{\titlerule[0.4pt]}]
%modify the format of the subsection headers, removes numeration
\titleformat{\subsection}{\normalfont\normalsize\bfseries}{}{0pt}{}
\setlength{\headheight}{80pt} % Increase the size of the header to accommodate meta-information
\pagestyle{fancy}\fancyhf{} % Use the custom header specified below
\renewcommand{\headrulewidth}{0pt} % Remove the default horizontal rule under the header 

%----------------------------------------------------------------------------------------
%	HEADER INFORMATION
%----------------------------------------------------------------------------------------
\def\proyecto{ } %nombre del proyecto de dise�o de m�quinas
\def\calculo{ } %nombre del c�lculo presentado
\def\nombre{ } % nombre de la persona que realiza los c�lculos
\def\revisor{ } %nombre de quien revisa los c�lculos
\def\fecha{ } % fecha de realizaci�n de los c�lculos

%----------------------------------------------------------------------------------------
%	HEADER CREATION
%----------------------------------------------------------------------------------------
\fancyhead[L]{\begin{tabular}{| l l | l l |} % The header is a table with 4 columns
\hline
\textbf{Proyecto}   & \proyecto  & \textbf{P�gina}    & \thepage/\pageref{LastPage} \\ % Project name and page count
\textbf{C�lculo}    & \calculo   & \textbf{Fecha}     & \fecha\\ % Fecha en la que se raliza el c�lculo
\textbf{Calculista} & \nombre    & \textbf{Revisor}   & \revisor \\ % Noombre del calculista y el revisor
\hline
\end{tabular}}

%------------------------------------------------------------------------------%
\begin{document}
%------------------------------------------------------------------------------%
%------------------------------------------------------------------------------%
\section{Descripci�n}\label{sec:descripci�n}




%------------------------------------------------------------------------------%
%------------------------------------------------------------------------------%
\section{Objetivo del c�lculo}\label{sec:objetivo}




%------------------------------------------------------------------------------%
%------------------------------------------------------------------------------%
\section{Variables del modelo de c�lculo}\label{sec:variables}



%\clearpage
%------------------------------------------------------------------------------%
%------------------------------------------------------------------------------%
\section{Imagen del modelo de c�lculo}\label{sec:imagenes}





%\FloatBarrier
%\clearpage
%------------------------------------------------------------------------------%
%------------------------------------------------------------------------------%
\section{Valores del modelo de c�lculo}\label{sec:valores}





%\FloatBarrier
%\clearpage
%------------------------------------------------------------------------------%
%------------------------------------------------------------------------------%
\section{Modelo de c�lculo}\label{sec:modelo}




%------------------------------------------------------------------------------%
%------------------------------------------------------------------------------%
\section{Resultados y conclusiones del modelo de c�lculo}\label{sec:conclusion}



 

\end{document} 














